\documentclass[11pt,letterpaper]{article}
\usepackage{multicol}
\usepackage{calc}
\usepackage{ifthen}
\usepackage[margin=0.4in]{geometry}
\usepackage{amsmath,amsthm,amsfonts,amssymb}
\usepackage{color,graphicx,overpic}
\usepackage{hyperref}
\usepackage{enumitem}

% Colors for design
\definecolor{sectioncolor}{RGB}{139, 69, 19} % Warm brown for Art
\definecolor{highlight}{RGB}{255, 248, 240}

% Custom section commands
\newcommand{\mysection}[1]{%
    \vspace{-0.3em}
    \colorbox{sectioncolor}{\makebox[\linewidth][l]{\bfseries\color{white}#1}}%
    \vspace{0.3em}
}

\newcommand{\mysubsection}[1]{%
    \vspace{-0.2em}
    {\bfseries\color{sectioncolor}#1}%
    \vspace{0.1em}
}

% Basic layout
\pagestyle{empty}
\makeatletter
\renewcommand{\section}{\@startsection{section}{1}{0mm}{-1ex plus -.5ex minus -.2ex}{0.5ex plus .2ex}{\normalfont\large\bfseries}}
\renewcommand{\subsection}{\@startsection{subsection}{2}{0mm}{-1explus -.5ex minus -.2ex}{0.5ex plus .2ex}{\normalfont\normalsize\bfseries}}
\makeatother

\setcounter{secnumdepth}{0}
\setlength{\parindent}{0pt}
\setlength{\parskip}{0pt plus 0.5ex}

\begin{document}
\raggedright
\footnotesize
\begin{multicols*}{2}

% ================= HEADER =================
\begin{center}
     \Large{\textbf{\color{sectioncolor}Art Appreciation Cheat Sheet}} \\
     \small{Prelims: Power of Art \& Everyday Aesthetics}
\end{center}

% ================= PART 1: POWER OF ART =================
\mysection{Part 1: The Power of Art}

\mysubsection{Learning How to See}
\begin{itemize}[leftmargin=*, noitemsep, topsep=0pt]
    \item Art requires \textbf{active viewing}, not passive looking
    \item Artists see details and relationships others miss
    \item Modern humans are visually numb from image overload
\end{itemize}

\mysubsection{The Mona Lisa (Leonardo da Vinci)}
\begin{itemize}[leftmargin=*, noitemsep, topsep=0pt]
    \item \textbf{Sfumato}: Soft, hazy transitions around eyes/mouth
    \item Mismatched horizon lines create illusion of movement
    \item Stolen 1911: the theft made it world-famous
    \item \textbf{Sprezzatura}: Aristocratic aloofness, effortless elegance
\end{itemize}

\mysubsection{Japanese Art Classification}
\begin{tabular}{l|l}
    National Treasures & Can't sell or export \\
    Important Cultural Properties & Significant but lower tier \\
    Living National Treasures & Title for living masters \\
\end{tabular}

\mysubsection{Art and Faith}
\begin{itemize}[leftmargin=*, noitemsep, topsep=0pt]
    \item \textbf{Venus of Willendorf}: Fertility charm (25,000 BCE)
    \item \textbf{Lascaux Cave Paintings}: Hunting magic, not decoration
    \item \textbf{African Masks}: Become the spirit when worn in ritual
    \item \textbf{Shamanism}: Everything has a spirit, art helps travel between worlds
    \item \textbf{Gothic Cathedrals}: Vertical lines pull soul upward
\end{itemize}

\mysubsection{Greek Idealism}
\begin{itemize}[leftmargin=*, noitemsep, topsep=0pt]
    \item Shows humans as they \textit{should} be, not as they are
    \item Perfect balance between body and mind
    \item ``Beauty is truth, truth beauty'' (John Keats)
\end{itemize}

\mysubsection{Art as Power Declaration}
\begin{itemize}[leftmargin=*, noitemsep, topsep=0pt]
    \item \textbf{Henry VIII (Holbein)}: Wide stance, direct gaze = intimidation
    \item \textbf{Mussolini Head (Bertelli)}: 360° rotation = sees everything
    \item \textbf{Futurism}: Art movement about speed, machines, modernity
    \item \textbf{Mughal Art}: Emperor with halo = divine right
\end{itemize}

\mysubsection{Art as Shock / Social Commentary}
\begin{itemize}[leftmargin=*, noitemsep, topsep=0pt]
    \item \textbf{Aunt Jemima (Betye Saar)}: Transforms racist symbol into resistance
    \item \textbf{L.H.O.O.Q. (Duchamp)}: Mona Lisa with mustache = Dadaism
    \item \textbf{Dadaism}: Anti-art, mocks tradition and sacred art
\end{itemize}

\mysubsection{Art and Emotion}
\begin{itemize}[leftmargin=*, noitemsep, topsep=0pt]
    \item \textbf{Vietnam Memorial (Maya Lin)}: V-shaped black granite wall
    \item Polished surface reflects viewer: living face the dead
    \item Names chronological, not alphabetical: forces you to search
\end{itemize}

\mysubsection{Transforming the Ordinary}
\begin{itemize}[leftmargin=*, noitemsep, topsep=0pt]
    \item \textbf{Urban Landscape (Zhan Wang)}: City made of kitchenware
    \item \textbf{The Gates (Christo)}: 7,503 saffron gates in Central Park
    \item \textbf{House (Whiteread)}: Cast of empty space inside building
\end{itemize}

\mysubsection{Why Artists Create}
\begin{itemize}[leftmargin=*, noitemsep, topsep=0pt]
    \item \textbf{Self-expression}: Frida Kahlo (pain as subject)
    \item \textbf{Play}: Paul Klee (whimsy, doodles)
    \item \textbf{Memory}: Marc Chagall (dreamlike nostalgia)
\end{itemize}

% ================= PART 2: EVERYDAY AESTHETICS =================
\mysection{Part 2: Everyday Aesthetics}

\mysubsection{Historical Context}
\begin{itemize}[leftmargin=*, noitemsep, topsep=0pt]
    \item Early: Aesthetics covered everything
    \item 18th/19th C: Narrowed to Fine Arts only
    \item 20th C: Philosophy of Art, separate from daily life
    \item Late 20th C: Return to Everyday Aesthetics
\end{itemize}

\mysubsection{Wabi-Sabi (侘寂)}
Japanese aesthetic: acceptance of transience and imperfection
\begin{itemize}[leftmargin=*, noitemsep, topsep=0pt]
    \item \textbf{Wabi}: Less is more, rustic simplicity
    \item \textbf{Sabi}: Attentive melancholy, passage of time
    \item Beauty in the \textbf{Imperfect, Impermanent, Incomplete}
\end{itemize}

\mysubsection{Kintsugi (金継ぎ)}
\begin{itemize}[leftmargin=*, noitemsep, topsep=0pt]
    \item Broken pottery repaired with \textbf{gold lacquer}
    \item Highlights cracks instead of hiding them
    \item Breakage becomes part of object's history
    \item Repaired object is \textit{more} beautiful than original
\end{itemize}

\mysubsection{6 Dimensions of Everyday Aesthetics}

\textbf{1. Broadening Scope}
\begin{itemize}[leftmargin=*, noitemsep, topsep=0pt]
    \item Beyond museums: nature, functional objects, chores, routines
\end{itemize}

\textbf{2. Functional Beauty}
\begin{itemize}[leftmargin=*, noitemsep, topsep=0pt]
    \item Usefulness can \textit{make} something beautiful
    \item Well-designed tools ARE beautiful
\end{itemize}

\textbf{3. Active Participation vs. Passive Viewing}
\begin{itemize}[leftmargin=*, noitemsep, topsep=0pt]
    \item Traditional: Spectator (stand back)
    \item Everyday: Agent (doing, engaging)
    \item \textbf{Pedestal Problem}: Putting everyday object on display destroys its everydayness
\end{itemize}

\textbf{4. Social Aesthetics}
\begin{itemize}[leftmargin=*, noitemsep, topsep=0pt]
    \item Interactions have aesthetic dimensions
    \item Tone of voice, gestures, how you handle objects
\end{itemize}

\textbf{5. Beyond Beauty and Sublimity}
\begin{tabular}{l|l}
    The Boring & Background experience \\
    The Plain & Simple, quiet \\
    The Messy & Chaotic, lived-in \\
    The Cute & Small, endearing \\
    The Cozy & Warm, intimate \\
\end{tabular}

\textbf{6. Negative Aesthetic Qualities}
\begin{itemize}[leftmargin=*, noitemsep, topsep=0pt]
    \item \textbf{Artistic lens}: View ugly thing as if it were art
    \item \textbf{Motivated action}: Ugliness compels you to change it
\end{itemize}

\mysubsection{Graffiti: Art or Pollution?}
\begin{itemize}[leftmargin=*, noitemsep, topsep=0pt]
    \item \textbf{Visual pollution}: Vandalism, no consent, eyesore
    \item \textbf{Social commentary}: Political messages, voice for marginalized (e.g., Banksy, Lumad murals)
\end{itemize}

\mysubsection{Criticisms of Everyday Aesthetics}
\begin{itemize}[leftmargin=*, noitemsep, topsep=0pt]
    \item \textbf{Subjective relativism}: Too personal for critical discussion
    \item \textbf{Self-indulgence}: Becomes just ``what I like''
    \item \textbf{Loss of authority}: If everything is aesthetic, does expertise matter?
\end{itemize}

\mysubsection{Related Art Movements}
\begin{itemize}[leftmargin=*, noitemsep, topsep=0pt]
    \item \textbf{Conceptual Art}: Idea matters more than object
    \item \textbf{Environmental Art}: Merged with landscapes
    \item \textbf{Performance Art}: Body and actions are the art
    \item \textbf{Socially Engaged Art}: Collaboration with public
\end{itemize}

\end{multicols*}
\end{document}
